\section{The exact codimension-two closure and the required rank-two correction}
\label{sec:pole-free-closure}

The two cutoff moments are continuous for the inherited Hilbert
norms in the regimes of MA-MO-010. Their common kernel is therefore
a proper closed subspace of the usual theta ideal. We identify it,
including its actual Riesz Gram matrix, and prove that every nonzero
admissible smoothing still has a rank-two outgoing polar-moment
defect. Orthogonal removal of that defect produces a defined
endomorphism, but changes the original source action.

Retain all hypotheses, spaces and notation of MA-MO-010. In
particular \(0<a<1/2\), \(s\ge0\), and either the fixed-gap
smooth regime or the unconditional generic-width Gevrey regime
holds. Write \(J=J_{a,s}\), \(K=K_{a,s,\phi}\), and
\[
 w(x)=(e^{2ax}+e^{-2ax})(1+x^2)^{s/2},\qquad
 H_\Sigma=L^2(\mathbb R,w(x)\,dx).
 \tag{MA-MO.137}
\]
Inner products in this section are linear in the first variable.

\paragraph{MA-MO-011 (exact closure, noninvariance for every nonzero test,
and canonical correction).}
There are continuous linear functionals \(\lambda_0,\lambda_1\)
on \(H_\Sigma\) whose restrictions to the actual theta source are
\[
 \lambda_0(\Theta f)=\ell_0(f),\qquad
 \lambda_1(\Theta f)=\ell_1(f)
                         \qquad(f\in\mathscr U_0).
 \tag{MA-MO.138}
\]
On \(J\) they are independent. If
\(\Lambda=(\lambda_0,\lambda_1):J\to\mathbb C^2\), then
\[
                  K=J\cap\ker\lambda_0\cap\ker\lambda_1
                    =\ker\Lambda,
                    \qquad \operatorname{codim}_J K=2.
 \tag{MA-MO.139}
\]
Let \(r_0,r_1\in J\) be their Riesz vectors and
\(M=\Lambda\Lambda^*\). This positive definite two-by-two
matrix, with entries \(M_{ij}=\lambda_i(r_j)\), gives
\[
        P_K=1_J-\Lambda^*M^{-1}\Lambda.
 \tag{MA-MO.140}
\]

For every nonzero test \(g\) admitted in MA-MO-010, the actual
trace-class map \(B_{\phi,g}:K\to J\) satisfies
\[
             \Lambda B_{\phi,g}:K\longrightarrow\mathbb C^2
                       \quad\hbox{surjective}.
 \tag{MA-MO.141}
\]
In particular its range is not contained in \(K\). The corrected
map and its defect are
\[
 \begin{aligned}
 B_{\phi,g}^{\rm corr}&=P_KB_{\phi,g}:K\longrightarrow K,\\
 C_{\phi,g}&=\Lambda^*M^{-1}\Lambda B_{\phi,g}:K\longrightarrow J,
              &\operatorname{rank}C_{\phi,g}&=2.
 \end{aligned}
 \tag{MA-MO.142}
\]
The first map is trace class. The correction has the minimum
possible rank among linear corrections whose subtraction makes
the range lie in \(K\). Its construction uses the actual inherited
Hilbert Gram matrix, not the additive \(L^2\) norm.

\begin{proof}
\emph{Continuous moment functionals in the actual norm.}
Put \(c=1/2-a>0\) and
\[
 p_a(x)=e^{cx}\phi(x),\qquad
 \widehat p_a(t)=\mathcal F_+p_a(t),\qquad
 \beta_a(t)=j_a(t)\widehat p_a(-t),
             \qquad j_a(t)=\zeta(\tfrac12+a+it)^{-1}.
 \tag{MA-MO.143}
\]
The function \(p_a\) is Schwartz: it decays exponentially
at the left and vanishes sufficiently far right, with the
same properties for its derivatives. In the fixed-gap regime,
the polynomial reciprocal estimate therefore gives
\(\beta_a\in L^2\).

In the generic regime a sharper estimate is needed. Integration
by parts, with vanishing boundary terms at both ends, gives
\[
 \widehat p_a(t)=
       -\frac{\mathcal F_+(e^{c\,\cdot}\phi')(t)}{c+it}.
 \tag{MA-MO.144}
\]
The numerator is the transform of a compact Gevrey function
of the same finite order: the product rule bounds its
\(n\)-th derivative by \(C A^n(n!)^q\), since the exponential
factor is taken on a fixed compact interval. Repeated
integration by parts and optimization in \(n\), as in
MA-RS.19, yield a bound \(C\exp(-d|t|^{1/q})\).
The denominator has modulus at least \(c\). This decay beats
the proved bound \(\exp(C'\log^2(2+|t|))\) for \(j_a\).
Thus \(\beta_a\in L^2\) in this regime too. Formula
(MA-MO.144) verifies the noncompact weighted-step estimate
without treating \(p_a\) itself as compactly supported.

For \(y\in H_{a,0}\), define
\[
        \lambda_0(y)=\int_{\mathbb R}(V_ay)(t)\beta_a(t)\,dt,
                 \qquad \lambda_1(y)=\lambda_0(Ry)
                       \quad(y\in H_{-a,0}).
 \tag{MA-MO.145}
\]
Each integral is absolutely controlled by Cauchy--Schwarz.
Since \(s\ge0\), the inherited sum norm dominates both
component norms, so these functionals are continuous on
\(H_\Sigma\), with
\(\|\lambda_i\|\le\|\beta_a\|_2\).

For an actual source vector \(y=\Theta f\), write
\(F(x)=2e^{x/2}f(e^x)\) as in MA-MO.127.
The checked domain identity gives \(j_aV_ay=V_aF\).
The bilinear Plancherel identity, with the unitary Fourier
normalization used throughout, therefore gives
\[
 \begin{aligned}
 \lambda_0(y)
 &=\int (V_aF)(t)\widehat p_a(-t)\,dt\\
 &=\int e^{ax}F(x)p_a(x)\,dx
   =\int e^{x/2}\phi(x)F(x)\,dx
   =\int_{\mathbb R}v(t)f(t)\,dt.
 \end{aligned}
 \tag{MA-MO.146}
\]
Poisson gives \(Ry=\Theta\mathfrak Ff\), proving the
second identity in (MA-MO.138). In particular these are
functionals for the actual theta norm, not an assumed
bounded inverse into an additive source norm.

\emph{The exact kernel closure and its actual Gram matrix.}
Independence of \(\ell_0,\ell_1\) on \(\mathscr U_0\)
was proved in MA-MO-008. Thus \(\Lambda\) has rank two
already on \(\Theta\mathscr U_0\), and hence on \(J\).
Choose \(y_0,y_1\) in that smooth theta image with
\(\Lambda y_i=e_i\), the two standard basis vectors.
Continuity immediately gives \(K\subseteq\ker\Lambda\).
Conversely, if \(y\in J\) and \(\Lambda y=0\), choose
\(z_n\in\Theta\mathscr U_0\) tending to \(y\) in the
inherited norm and set
\[
             z_n'=z_n-\lambda_0(z_n)y_0-\lambda_1(z_n)y_1.
 \tag{MA-MO.147}
\]
Then \(z_n'\in\Theta\mathscr U_{\rm pf}\) and
\(z_n'\to y\). This proves (MA-MO.139) and its exact
codimension, without assuming that finite codimension
survives completion automatically.

Here is an explicit construction of the Riesz vectors in
the declared metric. Set
\[
 h_a=V_a^{-1}\overline{\beta_a},\qquad
 d_a(x)=\frac{e^{2ax}}{w(x)}h_a(x),\qquad
 r_0=P_Jd_a,\qquad r_1=Rr_0.
 \tag{MA-MO.148}
\]
The vector \(h_a\) belongs to \(H_{a,0}\), and
\[
 \|d_a\|_{H_\Sigma}^2
  =\int\frac{e^{4ax}}{w(x)}|h_a(x)|^2\,dx
  \le\int e^{2ax}|h_a(x)|^2\,dx.
\]
Direct substitution shows that \(d_a\) represents
\(\lambda_0\) on \(H_\Sigma\). Orthogonal projection
therefore gives its representative on \(J\). Reflection
is unitary in this even weight, preserves \(J\) by the
theta--Fourier identity, and commutes with \(P_J\).
Thus \(r_1=Rr_0\) represents \(\lambda_1\) on \(J\).
Their independence makes \(M=\Lambda\Lambda^*\)
positive definite, and the standard orthogonal decomposition
\(J=\ker\Lambda\oplus\operatorname{ran}\Lambda^*\)
proves (MA-MO.140). For example,
\(\Lambda\Lambda^*M^{-1}\Lambda y=\Lambda y\), so the
displayed complementary projection removes exactly both
moments, and its range is orthogonal to their common kernel.

\emph{The exact outgoing moments of the smoothed source action.}
For the fixed nonzero test \(g\), define even compact
smooth additive functions
\[
 W_g(t)=\int g(z)e^{z/2}v(e^zt)\,dz,
             \qquad U=W_{g^\vee},\qquad W=W_g,
             \qquad g^\vee(z)=g(-z).
 \tag{MA-MO.149}
\]
Compactness and smoothness follow from the bounded dilation
parameter and compact smooth \(v\); near zero the first
function is the constant \(\mathcal Mg(1/2)\).
Both \(W\) and \(U\) are nonzero. Indeed, if \(W_g=0\),
differentiate its restriction to positive arguments in
logarithmic coordinates and take the ordinary Fourier
transform. It gives
\[
      \mathcal Mg(1/2-it)P(t)=0\quad(t\in\mathbb R),
                  \qquad P(t)=\int e^{itx}\phi'(x)\,dx.
 \tag{MA-MO.150}
\]
Here \(P(0)=-1\). The entire function
\(\mathcal Mg(1/2-it)\) must vanish on an interval, hence
identically; Fourier uniqueness implies \(g=0\), a
contradiction. The same argument applies to \(g^\vee\).

Let \(\mathcal D_g h(t)=\int g(z)e^{-z/2}h(e^{-z}t)\,dz\),
the actual additive dilation operator in (MA-MO.136).
Changing variables in its bilinear pairing gives
\(\int v\mathcal D_g h=\int W_g h\).
Also \(\mathfrak F\mathcal D_g
=\mathcal D_{g^\vee}\mathfrak F\), directly by Fourier
scaling. The operator \(L\) is its own transpose for
this bilinear pairing and commutes with \(\mathfrak F\).
For every \(f\in\mathscr U_{\rm pf}\), therefore,
\[
 \begin{aligned}
 \lambda_0(B_{\phi,g}\Theta f)&=\int f(t)(LW)(t)\,dt,\\
 \lambda_1(B_{\phi,g}\Theta f)&=\int f(t)(\mathfrak F LU)(t)\,dt.
 \end{aligned}
 \tag{MA-MO.151}
\]
No adjoint conjugation has been silently introduced: all
these source moment identities are complex bilinear, and
Fourier Fubini supplies the stated transpose identities.

\emph{Independence of the two outgoing moments modulo all input constraints.}
We show that the five even smooth functions
\[
                 1,\quad v,\quad\mathfrak Fv,
                       \quad LW,\quad\mathfrak F LU
 \tag{MA-MO.152}
\]
are linearly independent on \(\mathbb R\setminus\{0\}\).
The first three are independent by MA-MO-008. Suppose
\[
        A LW+B\mathfrak F LU=c_0+c_1v+c_2\mathfrak Fv.
\]
All terms except \(c_0\) are Schwartz, so \(c_0=0\).
Write \(q=1-v\). Outside a sufficiently large compact
interval the remaining identity is
\[
           -A\mathfrak F(v\mathfrak FW)+B\mathfrak F(qU)
                       =c_2\mathfrak Fv.
\]
The preimages of all three transforms are compact smooth
functions. Their transforms are entire, so equality on
that exterior interval holds everywhere. Fourier inversion
in the even sector gives
\[
                  -A v\mathfrak FW+BqU=c_2v.
 \tag{MA-MO.153}
\]
Near zero \(v=1\) and \(q=0\). If \(A\ne0\), this
forces the entire function \(\mathfrak FW\) to be
constant on an interval, hence constant everywhere.
Its Schwartz decay then forces \(W=0\), contrary to
(MA-MO.150). Thus \(A=0\), and (MA-MO.153) near zero
also gives \(c_2=0\).

If \(B\ne0\), (MA-MO.153) now gives \(qU=0\).
Since
\(\mathfrak F LU=\mathfrak F(qU)-v\mathfrak FU\),
the original identity becomes
\(-Bv\mathfrak FU=c_1v\). On the same neighborhood
of zero, \(\mathfrak FU\) is constant. The preceding
entire-function argument forces \(U=0\), again a
contradiction. Hence \(B=0\), and the remaining coefficients
vanish by the independence of the first three functions.
This proves (MA-MO.152).

The independence is realized by actual source tests. Choose
five distinct positive points with invertible evaluation
matrix for (MA-MO.152); they exist because the evaluation
vectors span \(\mathbb C^5\). Replace each point and its
negative by an even compact smooth bump of mass one,
supported away from zero. For sufficiently narrow bumps
their moment matrix remains invertible, as in MA-MO.97.
Solving that matrix with right side
\((0,0,0,z_0,z_1)^{\mathsf T}\) gives an even compact
\(f\), supported away from zero, with
\(f\in\mathscr U_{\rm pf}\) and outgoing moments
\((z_0,z_1)\) in (MA-MO.151). Thus surjectivity in
(MA-MO.141) already holds on genuine smooth source
vectors, and passes to the stated Hilbert domain.

\emph{The correction and its necessary rank.}
Formula (MA-MO.140) gives (MA-MO.142), and composition
of the trace-class map \(B_{\phi,g}\) with this bounded
projection is trace class. Its removed part has range
\(\operatorname{ran}\Lambda^*\), which is two dimensional.
Surjectivity of \(\Lambda B_{\phi,g}\) makes that range
exactly two dimensional. If any linear correction \(C\)
has \((B_{\phi,g}-C)(K)\subseteq K\), then
\(\Lambda C=\Lambda B_{\phi,g}\), so \(\operatorname{rank}C\ge2\).
The orthogonal correction attains this lower bound.
\end{proof}

\paragraph{Trace and source scope.}
The corrected endomorphism has a legitimate ordinary nuclear
trace and Fredholm determinant. If desired, bounded cyclicity
between the two Hilbert spaces gives the exact identities
\[
 \begin{aligned}
 \operatorname{Tr}_K(P_KB_{\phi,g})
      &=\operatorname{Tr}_J(B_{\phi,g}P_K),\\
 \det_K(1-zP_KB_{\phi,g})
      &=\det_J(1-zB_{\phi,g}P_K).
 \end{aligned}
 \tag{MA-MO.154}
\]
These use the actual trace-class maps just proved, as in
the inherited nuclear-pairing/Sylvester argument of MA-GP-008.
They are not an evaluation by an arithmetic explicit formula.
The operator \(P_KB_{\phi,g}\) is a changed source action,
depending on the cutoff and the inherited weight through
the Gram matrix. It does not agree with the uncorrected
source action on all of \(K\); their defect has rank two
for every nonzero admissible \(g\), including nonzero
admissible tests satisfying \(g(0)=0\).

The proper closure (MA-MO.139) resolves the closure question
left open in MA-MO-010 in its two stated regimes. Nothing
here asserts the same functional bounds on every exceptional
width with arbitrary smooth cutoff. No new arithmetic
zero-location assertion or polynomial interpolation cost
bound is obtained. The finite polar representation and
the inherited cote~1 nuclear-trace locators, scans~079--080
and~111--112 via BR-TR-01/02, distinguish the finite-rank
correction from the analytic construction of the maps;
finite-length \(K_0\) data alone neither choose its Gram
metric nor identify its trace with the original formula.
